\documentclass[11pt]{article}
\usepackage[a4paper,margin=23mm]{geometry}
\usepackage{fontspec,amsmath,amssymb,booktabs,graphicx,tabularx,hyperref,placeins}
\setmainfont{chinese.ttf}[Path=fonts/,BoldFont=chinese.ttf,BoldFeatures={FakeBold=1.5},ItalicFont=chinese.ttf,ItalicFeatures={FakeSlant=0.15}]
\XeTeXlinebreaklocale "zh"
\XeTeXlinebreakskip=0pt plus 1pt minus 0.5pt
\emergencystretch=3em
\linespread{1.12}
\setlength{\parindent}{2em}
\hypersetup{colorlinks=true,linkcolor=blue,citecolor=blue,urlcolor=blue,pdftitle={催化数据的出版来源偏移下，少量目标标签学到了什么}}
\renewcommand{\abstractname}{摘要}
\renewcommand{\refname}{参考文献}
\renewcommand{\figurename}{图}
\renewcommand{\tablename}{表}
\newcommand{\pp}{\,pp}
\newcommand{\styunit}{mg\,MeOH\,h$^{-1}$\,g$_{cat}^{-1}$}
\newcommand{\plotfigure}[3]{\begin{figure}[htbp]\centering\includegraphics[width=\linewidth]{figures/#1.pdf}\caption{#2}\label{#3}\end{figure}}
\newcommand{\NewScalarCount}{10,948,000}
\newcommand{\HistoricalCells}{42,525}
\newcommand{\OpenedCells}{5,600}
\newcommand{\RbfMatchMae}{2.828}
\newcommand{\RobustMae}{2.913}

\title{催化数据的出版来源偏移下，\\少量目标标签学到了什么}
\author{}\date{}
\begin{document}\maketitle
\begin{center}\small 第11步英文主文的完整中文译文\end{center}
\begin{abstract}
来自目标出版来源的少量测量可以改善催化预测，但这种改善本身并不能说明究竟迁移了哪些信息。我们在甲醇合成和甲烷氧化偶联（OCM）的来源隔离研究中，区分响应水平、输入与响应的对应关系、历史响应结构，以及测量位置。四种响应和两种OCM表征构成五个分析面板。一个由支持集决定的投影恒等式区分了目标本地拟合与历史响应复用；随后，通过排列支持标签，在保持均值不变的条件下检验对应关系。在五个标签的预算下，已有的面向查询集的V最优设计，使参数保持不变的混合预测器在各来源等权的平均绝对误差上取得改善：甲醇生产率按原生单位计，从62.608降至50.274；加入组成信息的OCM从3.081降至2.708个百分点。两项改善均通过了此前五面板主要比较家族的多重性校正。化学信息还可以改善只使用操作条件的预测器的选样。冻结管线在91个额外文献出版来源上的应用支持了校准的价值，而1739条由原研究作者报告的测量表明，较强的支持集本地RBF模型可以优于历史响应混合。在该出版来源的二十次催化剂身份分离划分中，基于组成的选样使RBF平均误差从4.922降至4.107，并在十九次划分中取得改善。预测算子风险表达式及其精确的谱协方差敏感性界，解释了为什么名义上具有信息量的测量位置不一定使实际误差最小。由此形成的基准识别了标签何时提供来源水平、本地响应关系或有用的测量几何，并给出相应的迁移主张边界。
\end{abstract}

\noindent\textit{译文说明：正文、图注和表格说明均为完整翻译；公式、实验数值、英文原图和参考文献保持原稿。支持集指提供已知实验结果的样本，查询集指待预测并在评估时计分的样本。}

\section{引言}
文献催化数据汇集了不同出版来源中的材料、操作条件、反应器和测量实践。这种异质性带来了广泛的实验覆盖，也使实际部署的目标变得至关重要。预测一个已在训练数据中出现的出版来源的另一条记录，与预测一个新的出版来源，是不同的任务：相关训练记录可能传递该来源的响应水平和特征性趋势。因此，在估计模型对未见来源的表现时，需要采用分组验证\cite{roberts}。当获得新的目标来源测量时，这一区别同样成立。测量的价值取决于测了什么，以及这些响应如何进入预测器。

近期的甲醇和OCM数据资源使上述问题可以在实测数据上开展研究\cite{toldy,yan,suvarna,schmack,mine}。由于催化测量需要样品制备、反应器时间和分析工作，少量目标标签的预算具有现实吸引力。校准迁移为调整响应关系提供了成熟工具\cite{feudale}，而回归混合可以稳定本地响应与历史修正之间的选择\cite{breiman,efron}。然而，相对于未适配历史模型取得较大改善，并不足以证明其学到的响应结构成功迁移。输出均值，或在同一支持标签上拟合的简单回归器，可能已经提供了大部分改善。

测量位置带来了第二个问题。实验设计方法引导测量走向具有信息量的输入\cite{yubi,gorodetsky2016}，其收益取决于使用这些标签的预测器，以及需要预测的查询样本。更详细的化学信息，即使不进入部署预测器，也可能改善测量选择。相反，适合普通插值的选样规则，在所有待查询催化剂身份都被排除出支持池时，可能表现不同。近期关于主动学习的分析强调了这种对工作流程的依赖\cite{nair2026,tang2026}。我们通过实测催化数据中受控的响应信息和几何信息干预来研究这一点。

本文的贡献是一个分解信息作用、并将来源隔离的评估框架，包含三个相互关联的部分。第一，使用匹配的支持集本地预测与残差预测，揭示支持集拟合未能重建的历史响应成分。第二，通过标签排列以及选样与预测表征的交叉组合，区分对应关系与响应水平，并区分选样几何与预测输入。第三，在比较测量位置时，使用全因子设计固定标签预算、查询总体和预测器。针对实际预测算子的工作风险计算，以及一个协方差敏感性界，为名义选样风险的局限提供了可检验的解释。贡献在于整合的机制分析及其经验证据范围；岭回归、核回归、混合和V最优设计均为已有方法。

\section{数据与评估}
\subsection{反应系统和来源身份}
历史语料包含1245条配对的甲醇转化率与选择性记录，来自82个出版来源和80个泄漏组；1230条甲醇时空收率（STY）记录，来自109个来源；以及1794条OCM记录，来自408个出版条目和405个泄漏组。甲醇转化率与选择性结合了TheMeCat和一项甲醇回溯研究\cite{toldy,yan}；STY采用Suvarna数据集\cite{suvarna}；OCM采用Schmack补充数据\cite{schmack}。转化率、选择性和C$_2$收率的误差使用百分点，STY使用\styunit。这里有两个反应系统和四种响应。OCM仅操作条件（F0）与加入组成信息（F1）是同一批记录的两种配对表征，因此得到的是五个面板，而不是五个独立反应数据集。二十个甲醇出版来源身份同时出现在两套响应数据中。

来源指一个出版条目，并不等同于实验室或物理批次。重复引用和已识别的上游数据重叠被分配到泄漏组；外层划分时保持这些组完整。甲醇操作输入包括温度、压力、H$_2$/CO$_2$、气体时空速（GHSV）和BET比表面积。OCM操作输入包括温度、CH$_4$与O$_2$分压、二者比值、总压力和接触时间。F1增加了与排列顺序无关、按组分角色区分的组成描述符、已知摩尔量和用量缺失指示。化学词表、未见物种通道、缺失值填补和缩放，均使用被允许的历史训练数据拟合。组成编码保留已报告的组分角色和缺失状态，并不重建制备过程或表面结构。补充材料给出了记录行身份、输入质量、单位和来源。

\subsection{历史划分、标签和固定评估总体}
采用五个泄漏组折，并使用随机种子42、101、202、303和404重复。全局骨干模型在外层训练分区内部，依据三折分组折外预测的来源等权MAE选择。候选包括$\alpha\in\{0.1,1,10,100\}$的岭回归，以及使用100或300棵深度为2/3的树，或500棵深度为8的树的梯度提升。适配器正则化和混合先验将历史来源作为伪目标，在每个符合条件的来源和预算下进行十次不查看标签的支持抽取。外层查询标签不参与模型选择或预处理。一个单独的划分审计在记录随机划分与来源隔离划分下重新拟合固定设置的岭回归和提升模型，从而将划分选择与调参分开。

对于具有$n_s$行记录的留出来源，按稳定标识符进行的打乱将$\max(5,\lceil n_s/4\rceil)$行分配到查询集，其余进入支持池。完整预算曲线要求池中至少有十行。预算$k\in\{1,2,3,5,10\}$保持相同的来源总体：两种配对甲醇响应各有32个来源，STY有28个，OCM有22个。五十个随机支持顺序提供嵌套前缀，确定性选样也使用嵌套前缀。每项比较共享查询集；同一选样方式下不同预测器的比较使用相同的实际标签。不同全因子分析保留各自归档的支持抽取种子，因此随机方法的点估计可能略有不同；配对主张始终使用相应全因子分析内部的匹配对照。选样器知道查询输入，但查询响应直到计分时才可用。

次要的高温查询使用完整的高温区块，其中每个支持样本温度都低于每个查询样本温度。符合条件的来源数分别为：配对甲醇响应23个、STY 21个、历史OCM五个。只有五个来源的OCM压力测试属于探索性分析。补充材料还报告了甲醇催化剂名称分层留出；已报告的催化剂名称并不能证明出现了新的物理催化剂。

\subsection{额外文献及实测出版来源上的应用}
归档仓库提交版本的Mine OCM更新\cite{mine}含4759条记录，来自533个出版编号。计分之前，将完整历史参考文献表作为依据进行文献筛查，并删除三个重新列入的旧出版来源。符合要求的热催化OCM记录需具有可归属的额外引用、有效实测收率、一致单位和受支持的组成模式。每个排除项均保留原因。最终纳入2381条记录，来自91个额外出版来源，其中52个符合完整随机预算曲线条件，40个符合高温查询条件。这证明了在一个单独发布的汇编中出版来源不重叠，但并未排除共享实验室或上游抽取错误。历史拟合和调参均先于目标支持与查询计分。

第二项应用使用Nishimura等人报告的验证测量\cite{nishimura2023new}。原始2928行CADS文件混合了不同研究阶段。删除1024条阶段不符合条件的记录、120条缺少必要温度或响应的记录，以及45条重复的催化剂与条件记录后，得到1739条测量，覆盖116个催化剂标识符。它们属于一个出版来源和一个实验家族，并不会因为来自五个验证阶段就成为五个独立数据集。四项操作输入和组成用量不可用，在冻结的历史预处理下保持缺失，不猜测物理值进行填充。首次评估该出版来源时采用仅依靠历史数据冻结的管线。之后的风险方法比较和重复催化剂划分使用了已经打开的出版来源，属于开发阶段的敏感性分析。

\subsection{估计对象与不确定性}
对固定支持抽取，先在查询样本内平均MAE和MSE，再在每个来源内对抽取及重复种子平均；各来源权重相同。因此，拥有较多查询记录或随机支持抽取的来源不会主导结果。配对对比使用来源层面的差值；自助重采样保持泄漏组完整，并以已拟合历史模型为条件。此前五面板主要选样比较家族使用校正后的区间。后续分析保持原有的开发或探索性状态。重复折和重复催化剂划分均不增加独立出版来源数量。对于单一实测实验家族，我们报告点估计和条件性划分分布，不报告出版来源总体的置信区间。

\section{方法：分解目标标签中的信息}
\subsection{本地拟合与历史响应复用}
设$f_G$为冻结的历史预测器，$S$为有标签的支持集，$Q$为与其不相交的查询集，$z(x)$为由训练数据拟合的特征。本地岭回归拟合支持集响应：
\begin{equation}
q_L(x)=b_L+z(x)^\top\beta_L,\qquad
(b_L,\beta_L)=\arg\min_{b,\beta}\sum_{i\in S}(y_i-b-z_i^\top\beta)^2+\alpha\|\beta\|_2^2.
\end{equation}
残差岭回归则拟合$y_S-f_G(S)$，再加回$f_G(Q)$。当正则化相同且截距不受惩罚时，将响应映射到查询预测的岭回归算子$R_{S,\alpha}$是线性的，因此：
\begin{equation}\label{eq:projection}
q_R(Q)-q_L(Q)=f_G(Q)-R_{S,\alpha}[f_G(S)](Q).
\end{equation}
右侧是由支持集条件决定的拟合未能重建的历史响应结构。这并不保证误差改善。对本地和残差正则化分别调参，还会改变算子，所以补充材料包含匹配$\alpha$的消融。当只有一个标签时，本地岭回归等于该支持响应，残差岭回归则是全局预测加上一个残差偏移。

收缩混合为$q_M=(1-w)q_L+wq_R$。设$l_{-i}$和$r_{-i}$为支持集留一预测，$d_i=r_{-i}-l_{-i}$。历史先验提供$w_0$与分歧尺度$D^2$；仅使用支持集的估计为：
\begin{equation}\label{eq:blend}
\widehat w=\operatorname{clip}_{[0,1]}
\frac{\sum_i d_i(y_i-l_{-i})+cD^2w_0}{\sum_i d_i^2+cD^2}.
\end{equation}
先验网格为$w_0\in\{0,0.25,0.5,0.75,1\}$和$c\in\{\infty,10,2,0.5,0\}$。为伪来源选择训练先验时，排除其所属泄漏组。只有一个标签或退化情形下，使用相关冻结先验。其他对照包括未适配全局预测器、支持均值、残差偏移、带惩罚仿射校正、经验贝叶斯偏移，以及仅使用温度的岭回归。较强的本地RBF对照使用由历史数据选择的超参数。目标查询响应不进入任何预测接口。

\subsection{标签对应关系与选样表征}
在固定支持集内部排列响应，保留响应的多重集合及均值，但打破其与输入行的对应。因此，仅使用均值的预测不变。若依赖特征的预测变差，就能在给定支持集和历史拟合的条件下识别对应关系的价值。排列是一种信息干预，并不是关于化学组成对实测响应具有因果作用的主张。

选样采用面向查询集的线性V最优设计，在固定的有限截距先验和冻结正则化下，贪心地最小化指定查询处的后验方差。标准对照包括随机选样、特征覆盖、面向池的V设计、核V设计，以及在操作条件和组成两种视图间取稳健最大值的方法。表征交叉时，将仅操作条件选样和加入组成信息的选样分别与F0、F1预测配对，在两个预测器之间固定每个实际选出的支持集。仅在选样环节打乱化学信息的行对应关系，提供额外对照。完整矩阵和策略特定调参的比较见补充材料；主要测量位置对比保持预测器参数固定。

\subsection{预测算子风险与协方差敏感性}
对于固定的自由截距核预测器，设$K_{SS}$和$K_{QS}$为支持核矩阵和查询与支持之间的核矩阵，$C=K_{SS}+\alpha I$。其响应权重为：
\begin{equation}\label{eq:weights}
A=K_{QS}C^{-1}+
\frac{(\mathbf1-K_{QS}C^{-1}\mathbf1)\mathbf1^\top C^{-1}}
{\mathbf1^\top C^{-1}\mathbf1},\qquad \widehat y_Q=Ay_S.
\end{equation}
$A$每行的和均为一，因此常量形式的目标来源响应偏移会从预测误差中抵消。这是已有的中心化核岭回归或泛克里金方法\cite{gpml}。核候选库包含各向同性RBF，以及操作条件块和组成块的加性、乘积和ANOVA组合。历史随机支持验证将部署核固定在所有选样规则之间。采用相同调参过程的各向同性对照，将扩展核家族的作用与调参数据改变的作用分开。

令$B=[-A\ I]$为作用于堆叠支持与查询响应的误差算子，$\Sigma_0$为观测响应的工作协方差。当常量均值被正确表示时：
\begin{equation}\label{eq:risk}
\mathcal R_0(S)=|Q|^{-1}\operatorname{tr}(B\Sigma_0B^\top).
\end{equation}
当核的对角线为一、独立工作噪声为$\nu$时，它等于$1+\nu-2(A\overline K_{SQ})_{qq}+[A(\overline K_{SS}+\nu I)A^\top]_{qq}$的均值。协方差可以使用部署核，也可以使用由历史数据给出的核假设经验平均。平方根版本是高斯绝对损失的代理量。对协方差取平均并不构成经校准的后验，也不会给出高斯混合的精确绝对风险。

当协方差误差满足$\|\Sigma-\Sigma_0\|_2\leq\varepsilon$时，固定算子的敏感性为：
\begin{equation}\label{eq:bound}
|\mathcal R_\Sigma(S)-\mathcal R_0(S)|\leq
\varepsilon\left(1+\frac{\|A\|_F^2}{|Q|}\right).
\end{equation}
进一步地，在满足$\Sigma\succeq0$的这个谱范数球上，最坏风险恰好等于名义风险加上右侧项，由$\Sigma=\Sigma_0+\varepsilon I$达到。这个基本的稳健设计恒等式明确体现了支持噪声的放大。非常量均值失配会另外增加一个平方偏差项。该界假定给定了协方差半径，并没有证明真实催化数据的协方差一定处于该半径内。

我们使用同一个较强的、在五标签历史数据上选择的RBF，评估名义RBF匹配选样，以及最小化该惩罚风险的选样。惩罚$\varepsilon\in\{0,0.1,1,10\}$仅依据历史伪来源的来源/组等权五标签MAE选择。取零时提供精确的名义对照。五种预算均使用固定的RBF参数化；在五标签时，其与此前各预算单独调参的对照相同。完整十四策略历史比较保留了此前有限截距、中心化平方/绝对风险和协方差平均选样。

\subsection{重复排除催化剂身份}
对于实测出版来源，二十次稳定且不查看标签的打乱，每次将116个催化剂标识符中的29个分配到查询集，另87个分配到支持池。同一身份的全部记录始终处于同一侧。每次划分中，随机选样使用五十个嵌套顺序；操作V、组成V、池V、扩展核风险和两种RBF匹配规则使用各自完整的确定性顺序。冻结历史参数在各划分间共享，所有划分均被报告。这些重复衡量的是一个实验家族内对查询组成的敏感性，并不是二十个不同来源上的重复验证。

\section{结果}
\subsection{目标标签提供的不只是来源响应水平}
受控划分审计表明，记录随机划分表现与来源隔离表现回答了明显不同的预测问题（图\ref{fig:split}）。固定模型配置，使这种差异能够归因于验证设计，而非选择了更有利的模型。这支持在评估标签辅助部署时使用来源隔离。
\plotfigure{fig1_split_audit}{固定模型设置和训练拟合预处理的受控划分审计。按响应和模型家族分别报告记录随机与泄漏组隔离的误差。百分比响应和STY保持各自原生单位。本图重现已审计的历史拟合。}{fig:split}

在五个随机标签下，支持均值在所有五个历史面板中均优于全局预测器；依赖特征的本地拟合还提供了进一步改善（表\ref{tab:information}）。对于加入组成信息的OCM，MAE由未适配模型的4.527降至支持均值的3.693，并进一步降至本地岭回归的3.374。STY相应数值为139.258、91.993和65.210。这些比较将目标响应水平的价值与本地响应关系的价值分开。
\begin{table}[htbp]\centering\small
\caption{五标签历史信息对照，每个面板内部共享随机支持集。误差为来源等权MAE；全局对照不使用目标标签。STY单位为\styunit，其他响应单位为百分点。}\label{tab:information}
\begin{tabular}{lrrrrr}
\toprule
策略 & 转化率 & 选择性 & OCM F0 & OCM F1 & STY \\
\midrule
全局模型 & 5.271 & 22.740 & 5.302 & 4.527 & 139.258 \\
支持均值 & 4.304 & 17.136 & 3.693 & 3.693 & 91.993 \\
偏移 & 3.455 & 15.747 & 3.672 & 3.423 & 88.265 \\
仿射校正 & 3.083 & 12.539 & 3.694 & 3.360 & 71.225 \\
经验贝叶斯偏移 & 3.457 & 15.988 & 3.759 & 3.415 & 86.510 \\
本地岭回归 & 2.429 & 11.547 & 3.655 & 3.374 & 65.210 \\
温度岭回归 & 2.519 & 11.535 & 3.677 & 3.677 & 87.098 \\
残差岭回归 & 2.720 & 12.077 & 3.701 & 3.109 & 69.889 \\
收缩混合 & 2.441 & 11.584 & 3.660 & 3.081 & 62.608 \\
\bottomrule
\end{tabular}

\end{table}

排列提供了直接的对应关系诊断。在面向查询的V选样下，破坏输入与响应对应，使混合器MAE在转化率上增加2.015个百分点，选择性上增加6.040个百分点，STY上增加30.671个原生单位，OCM F1上增加0.840个百分点（表\ref{tab:alignment}）。这四个对比的名义分组自助区间均排除零。OCM F0仅增加0.047个百分点，区间包含零。干预保持相同的支持均值：标签能够提供仅凭边际响应水平无法获得的关系，而这种关系能够被检测到的价值依赖表征。
\begin{table}[htbp]\centering\small
\caption{相同五标签查询V设计下的支持标签排列。区间为探索性的条件95\%分组自助区间，不构成新的比较家族校正后主张。}\label{tab:alignment}
\begin{tabular}{lrrrl}\toprule
面板 & 正确对应 & 排列后 & $\Delta$ MAE & 95\%区间 \\ \midrule
转化率 & 2.121 & 4.135 & 2.015 & [1.253, 2.851] \\
选择性 & 10.683 & 16.723 & 6.040 & [3.693, 8.513] \\
STY & 50.274 & 80.945 & 30.671 & [10.436, 59.603] \\
OCM F0 & 3.768 & 3.815 & 0.047 & [-0.009, 0.134] \\
OCM F1 & 2.708 & 3.548 & 0.840 & [0.478, 1.237] \\
\bottomrule\end{tabular}

\end{table}

残差拟合和混合拟合并不会普遍优于本地拟合。历史F1混合器优于本地岭回归，而其他若干面板的改善很小或表现变差。匹配$\alpha$的对比隔离了式\ref{eq:projection}的成分，其区间在每个面板中均包含零。因此，获得标签带来的显著收益，可以与历史响应函数有限的增量价值同时存在。

\subsection{额外来源仍显示校准价值}
在91个出版来源的扩展中，固定随机查询完整曲线总体有52个来源。在五标签下，F1全局MAE为6.120个百分点，支持均值为4.490，本地岭回归为3.593，混合器为3.552（表\ref{tab:additional}）。这验证了在额外出版来源中有用的来源条件校准。混合器相对于本地模型的小幅增量，并没有验证历史数据上更大的增量。原始对比状态及其条件区间保留于补充材料；不会将该额外来源应用转化为普遍迁移主张。
\begin{table}[htbp]\centering\small
\caption{额外文献出版来源上的冻结应用，使用五个随机支持标签。随机查询有52个来源，高温查询有40个来源。原生MAE单位为百分点。}\label{tab:additional}
\begin{tabular}{lrrrr}
\toprule
策略 & F0随机 & F1随机 & F0高温 & F1高温 \\
\midrule
全局模型 & 5.973 & 6.120 & 5.699 & 5.140 \\
支持均值 & 4.490 & 4.490 & 5.556 & 5.556 \\
偏移 & 4.502 & 4.606 & 5.512 & 5.397 \\
仿射校正 & 4.284 & 4.410 & 5.209 & 5.328 \\
经验贝叶斯偏移 & 4.545 & 4.621 & 5.205 & 4.942 \\
本地岭回归 & 4.045 & 3.593 & 4.801 & 4.365 \\
温度岭回归 & 3.779 & 3.779 & 4.697 & 4.697 \\
残差岭回归 & 4.421 & 3.716 & 5.375 & 4.454 \\
收缩混合 & 4.054 & 3.552 & 4.806 & 4.335 \\
\bottomrule
\end{tabular}

\end{table}

实测出版来源应用通过原研究作者报告的验证数据体现了相同区别。在组成V支持集上，F1未适配模型的MAE为5.869，支持均值5.959，本地岭回归4.170，混合器4.213，历史调参的本地RBF为3.813。依赖特征的响应拟合有用，而历史响应混合并非最强预测器。操作字段和组成用量缺失限制了该应用；即使本地拟合不使用已拟合的历史响应函数，它在使用历史预处理时仍利用了历史输入知识。

\subsection{相同预算下，测量位置具有实际价值}
标准面向查询的V设计在若干历史面板上改善了参数保持不变的混合器（表\ref{tab:acquisition}）。STY从62.608降至50.274，点估计误差降低19.7\%；OCM F1从3.081降至2.708，降低12.1\%。在此前五面板主要比较家族中，两项校正后差值区间均排除零。转化率和选择性的点估计有利，但未通过该校正。OCM F0没有改善。因此，收益成立于指定表征和响应，而不是均匀成立于所有面板。
\begin{table}[htbp]\centering\small
\caption{此前五面板主要比较家族：相同冻结混合器和五标签预算下，查询V与随机选样比较。校正后区间针对原生单位MAE差值。}\label{tab:acquisition}
\begin{tabular}{lrrrrl}\toprule 面板 & $n$ & 随机 & 查询V & 降低(\%) & 校正后$\Delta$区间\\\midrule
转化率 & 32 & 2.441 & 2.121 & 13.1 & [-0.707, 0.068]\\
选择性 & 32 & 11.584 & 10.683 & 7.8 & [-2.140, 0.615]\\
STY & 28 & 62.608 & 50.274 & 19.7 & [-25.965, -2.874]\\
OCM F0 & 22 & 3.660 & 3.768 & -3.0 & [-0.355, 0.714]\\
OCM F1 & 22 & 3.081 & 2.708 & 12.1 & [-0.746, -0.056]\\
\bottomrule\end{tabular}

\end{table}
\plotfigure{step7_01_primary}{五个面板中相同预算查询V的增量。绘图时，原生单位校正后区间除以每个面板固定的实际随机MAE。它们不是自助法比值区间，也不是跨反应系统合并的效应。}{fig:primary}

选样化学信息与预测化学信息发挥不同作用。在22个历史OCM来源的匹配表征交叉中，组成V使仅操作条件混合器的MAE由随机支持下的3.705降至3.370；预测器从未接收组成输入。打乱选样化学信息后，误差上升至3.647。对F1，组成V与打乱化学信息后的相应误差分别为2.708和3.039。这些结果识别了独立于响应拟合特征、可用于选择测量的几何信息（图\ref{fig:crossing}）。
\plotfigure{step9_01_matched_geometry}{选样与预测表征交叉。每个实际选出的支持集在F0和F1预测间共享。不同选样规则在同预算下使用不同支持集。保留历史与实测出版来源数值；更有利的选样视图依赖部署预测器和来源。}{fig:crossing}

在首次实测出版来源应用中，组成V使F1混合器MAE从随机选样的4.787降至4.213，点估计降低12.0\%。对F0，操作V优于组成V。学习得到的视图门控在历史任务中几乎总是回退到组成视图，未能证明优于固定对照。这些发现支持了选样表征机制，同时保留其部署依赖性。

\subsection{查询身份与工作风险限制选样收益}
初始催化剂留出将29个身份的435行作为查询，将其他87个身份的1304行作为池。使用相同较强本地RBF和五标签，扩展核匹配选择的MAE为3.689，随机为4.818，组成V为4.680。其23.4\%点估计改善以该划分为条件。在普通随机查询中，同样的扩展风险选择却使RBF MAE恶化至5.145，随机为4.792，组成V为3.813。因此，即便预算与预测器设置相同，当查询任务改变时，选样方法的排序仍可能反转。

新增的二十次身份分离划分，概括了这一发现对查询组成的敏感程度（表\ref{tab:repeated}、图\ref{fig:sensitivity}）。每个随机对照均在同一划分中对五十个实际支持顺序取平均。显示范围是划分之间的条件性散布，并不是对未见出版来源的不确定性区间。
\begin{table}[htbp]\centering\small
\caption{已打开实测出版来源上的二十次新增催化剂身份分离划分。所有结果使用相同的固定历史选择RBF和五标签。$\Delta$为方法MAE减去同划分匹配随机MAE，单位为百分点。}\label{tab:repeated}
\begin{tabular}{lrrll}\toprule
选样 & 平均MAE & 平均$\Delta$ & 划分$\Delta$范围 & 改善次数 \\ \midrule
随机 & 4.922 & +0.000 & +0.000 至 +0.000 & 0/20 \\
组成V & 4.107 & -0.815 & -1.616 至 +0.120 & 19/20 \\
操作V & 4.926 & +0.004 & -1.025 至 +0.731 & 11/20 \\
池V & 4.221 & -0.701 & -1.272 至 +0.210 & 18/20 \\
扩展核风险 & 4.480 & -0.441 & -1.545 至 +0.498 & 15/20 \\
RBF匹配风险 & 4.652 & -0.270 & -1.360 至 +0.621 & 13/20 \\
惩罚RBF风险 & 4.652 & -0.270 & -1.360 至 +0.621 & 13/20 \\
\bottomrule\end{tabular}

\end{table}
组成V平均MAE为4.107，随机选样为4.922；以两个均值的比值计算，误差降低16.6\%。组成V在二十次划分中有十九次改善。池V改善十八次，扩展核风险改善十五次，RBF匹配风险改善十三次。因此，稳定的正向发现是在身份排除下，化学信息引导的测量覆盖具有价值；此前单个划分上的排序并不普遍成立。这些重复划分仍属于一个已打开实验家族内的敏感性分析。
\plotfigure{task_sensitivity}{选样依赖查询任务。A面板显示二十次催化剂划分相对于匹配随机对照的全部差值，箱线图概括其条件性分布。B面板保留此前普通查询与初始催化剂留出的结果。划分共享测量，属于一个实验家族。}{fig:sensitivity}

将选样匹配到较强各向同性RBF，得到历史五标签MAE \RbfMatchMae\ 个百分点；由历史数据选择的惩罚规则为\RobustMae{}（表\ref{tab:matched}）。历史伪来源调参在25个外层拟合中，有23个选择零惩罚，两个选择0.1；完整历史选择零。因此，惩罚未提供稳定的新优势。此前以协方差平均绝对风险为主要候选的方案，在家族校正后也未超过其五个较强对照。有限截距、中心化平方风险和中心化绝对风险选样器可能选出相同的五点支持集，因而它们共同取得的催化剂留出收益，无法识别某一种风险公式作为原因。这些对照说明，式\ref{eq:bound}适合作为敏感性的解释，而不是对优越算法的实证认证。
\begin{table}[htbp]\centering\small
\caption{固定较强RBF、五标签和22个符合条件来源下的历史随机查询F1比较。工作风险是无量纲代理量，并不是以百分点计的实际MAE估计。}\label{tab:matched}
\begin{tabular}{lrrr}\toprule
选样 & MAE(百分点) & 工作风险 & 放大系数 \\ \midrule
随机 & 3.288 & 0.137 & 0.563 \\
组成V & 2.851 & 0.091 & 0.583 \\
扩展核风险 & 2.784 & 0.090 & 0.627 \\
RBF匹配风险 & 2.828 & 0.090 & 0.664 \\
惩罚RBF风险 & 2.913 & 0.091 & 0.653 \\
\bottomrule\end{tabular}

\end{table}

\section{讨论}
目标标签改善的是条件问题$\widehat f(Q\mid X_Q,X_S,y_S,\mathcal D_H)$。它们揭示历史训练集$\mathcal D_H$中不存在的目标来源信息。一个响应可以校准水平，多个正确对应的响应可以刻画本地关系，经过设计选择位置的响应可以更好地覆盖查询样本。这些变化均不能证明对任意未来来源的零标签预测得到改善。少标签预算曲线也不要求每增加一个标签都下降：新测量响应可能改变带噪拟合，或改变不同修正方法之间的相对权重。

式\ref{eq:projection}中的区别具有实际意义。当适配方法主张复用了历史响应知识时，支持集本地响应拟合是必要对照。仅用于选样的表征仍可能具有价值，因为选择测量与使用响应进行预测是不同的信息路径。基于这些认识，在将收益归因于迁移前，应采用明确的来源身份、较强本地对照、共享支持标签，以及选样与预测交叉的实验。

风险匹配增加了第三个实际条件。最小化工作方差的意义，取决于指定的预测器、查询总体和协方差假设。对于保持常量的预测器，自由来源水平会抵消，但响应形状失配与支持放大仍然存在。谱界隔离了敏感性的一部分；真实协方差误差与非常量均值偏差尚未被认证。没有验证这些假设就增加方法复杂度，不能替代具体任务上的证据。重复催化剂划分和完整历史全因子设计使这种依赖关系可以检查。

本研究范围是对实测催化响应的回溯分析，并非新的湿实验验证或自主催化剂发现。文献来源不重叠并不能证明实验室、抽取流程或物理批次独立。历史OCM高温证据仅有五个符合条件来源。原研究作者实测应用中的缺失字段和出版选择限制了代表性。重复种子、来源记录数量、催化剂划分及数百万条标量预测，均不会扩大独立来源总体。自助区间以已拟合管线为条件，并未涵盖重新估计历史模型的全部不确定性。记录催化剂批次和反应器元数据、并预先指定选样的前瞻性测量计划，可以直接检验部署结论。

\section{结论}
少量目标标签具有可以区分的科学作用：估计来源响应水平、学习正确对应的本地响应关系，以及通过测量位置改善查询覆盖。来源隔离比较确立了有用的条件校准；配对信息干预则说明，为何不能自动将其收益归因于历史响应迁移。同预算选样在甲醇生产率和加入组成信息的OCM上产生了有证据支持的收益，表征交叉证明化学信息在选择测量上的价值超出了其作为预测输入的价值。查询身份敏感性和预测算子风险识别了这些选样收益发生变化的条件。发布的基准为选择测量和在证据支持的层面提出迁移主张提供了可复现基础。

\section*{作者信息与声明}
作者身份、单位、贡献说明、基金及利益冲突披露，由投稿作者在配套作者信息记录中填写。当前科学稿件为作者中性的投稿主稿。

\section*{数据可用性}
配套计算档案包含来源与记录行账本、分析代码、冻结模型参数和已拟合对象、机器可读指标、全部评估查询预测，以及复现说明。原始数据来源已在上文引用，其来源和再利用条款见补充材料。本研究未产生新的湿实验测量。持久化仓库记录及标识符应由作者安排期刊存证时提供；本稿不声称已经获得仓库DOI。

\begin{thebibliography}{99}\small
\bibitem{toldy} Toldy, \'A. I. et al. TheMeCat: Dataset of Thermocatalytic Conversion of CO$_2$ to Methanol. \emph{Scientific Data} \textbf{13}, 165 (2026). \url{https://doi.org/10.1038/s41597-025-06482-8}.
\bibitem{yan} Yan, M. et al. A machine learning perspective on three decades of methanol synthesis: research framework and experimental operation insights. \emph{Communications Engineering} \textbf{5}, 162 (2026). \url{https://doi.org/10.1038/s44172-026-00706-4}.
\bibitem{suvarna} Suvarna, M., Ara\'ujo, T. P. \& P\'erez-Ram\'irez, J. A generalized machine learning framework to predict the space-time yield of methanol from thermocatalytic CO$_2$ hydrogenation. \emph{Applied Catalysis B: Environmental} \textbf{315}, 121530 (2022). \url{https://doi.org/10.1016/j.apcatb.2022.121530}.
\bibitem{schmack} Schmack, R. et al. A meta-analysis of catalytic literature data reveals property--performance correlations for the OCM reaction. \emph{Nature Communications} \textbf{10}, 441 (2019). \url{https://doi.org/10.1038/s41467-019-08325-8}.
\bibitem{mine} Mine, S. et al. Analysis of Updated Literature Data up to 2019 on the Oxidative Coupling of Methane Using an Extrapolative Machine-Learning Method to Identify Novel Catalysts. \emph{ChemCatChem} \textbf{13}, 3636--3655 (2021). \url{https://doi.org/10.1002/cctc.202100495}. Data: \url{https://github.com/mts-uw/OCM}, commit \texttt{dbf7be5aadb30323f825e8f7d7f6bd3012c99dd2}.
\bibitem{roberts} Roberts, D. R. et al. Cross-validation strategies for data with temporal, spatial, hierarchical, or phylogenetic structure. \emph{Ecography} \textbf{40}, 913--929 (2017). \url{https://doi.org/10.1111/ecog.02881}.
\bibitem{feudale} Feudale, R. N. et al. Transfer of multivariate calibration models: a review. \emph{Chemometrics and Intelligent Laboratory Systems} \textbf{64}, 181--192 (2002). \url{https://doi.org/10.1016/S0169-7439(02)00085-0}.
\bibitem{efron} Efron, B. \& Morris, C. Stein's estimation rule and its competitors---an empirical Bayes approach. \emph{Journal of the American Statistical Association} \textbf{68}, 117--130 (1973).
\bibitem{breiman} Breiman, L. Stacked Regressions. \emph{Machine Learning} \textbf{24}, 49--64 (1996). \url{https://doi.org/10.1023/A:1018046112532}.
\bibitem{yubi} K. Yu, J. Bi and V. Tresp. Active learning via transductive experimental design. ICML, 1081--1088 (2006). \href{https://doi.org/10.1145/1143844.1143980}{doi:10.1145/1143844.1143980}.
\bibitem{gpml} C. E. Rasmussen and C. K. I. Williams. Gaussian Processes for Machine Learning. MIT Press (2006), Chapters 2 and 4. \url{https://gaussianprocess.org/gpml/}.
\bibitem{nishimura2023new} S. Nishimura, X. Li, J. Ohyama and K. Takahashi. Leveraging machine learning engineering to uncover insights into heterogeneous catalyst design for oxidative coupling of methane. Catal. Sci. Technol. 13, 4646--4655 (2023). \href{https://doi.org/10.1039/d3cy00596h}{doi:10.1039/d3cy00596h}. Author measurements: \url{https://cads.eng.hokudai.ac.jp/datamanagement/datasources/d84c1e22-ceb9-488a-8d45-4c7cf1c603b5/}.
\bibitem{gorodetsky2016} A. Gorodetsky and Y. Marzouk. Mercer Kernels and Integrated Variance Experimental Design: Connections Between Gaussian Process Regression and Polynomial Approximation. SIAM/ASA J. Uncertain. Quantif. 4, 796--828 (2016). \href{https://doi.org/10.1137/15M1017119}{doi:10.1137/15M1017119}.


\bibitem{tang2026} R. Tang, S. J. Sloman and S. Kaski. Representative, Informative, and De-Amplifying: Requirements for Robust Bayesian Active Learning under Model Misspecification. AISTATS, PMLR 300, 3016--3024 (2026). \url{https://proceedings.mlr.press/v300/tang26d.html}.

\bibitem{nair2026} A. S. Nair and L. Foppa. A critical examination of active learning workflows in materials science. \emph{Digital Discovery} \textbf{5}, 2366--2382 (2026). \url{https://doi.org/10.1039/d6dd00081a}.

\end{thebibliography}

\end{document}
